% À placer après \chapter{Conclusion générale}
\label{chap:conclusion}
%\rouge{20 mai - 1er juin}

Au cours de cette thèse, nous avons abordé et essayé d'explorer en détail le problème de clustering et de structuration automatique de données.
Nous sommes partis d'un cadre minimaliste, lorsque les données $\mathcal X \subset \mathbb R^p$ sont dans l'espace euclidien. Le modèle statistique le plus courant consiste à supposer que les données $\mathcal X = \bigl\{ x_1, \ldots , x_n \bigr\}$ sont le résultat de $n$ réalisations indépendantes et identiquement distribuées (i.i.d), selon une certaine loi à densité $f$ par rapport à la mesure de \nom{Lebesgue}.
Le but est de retrouver la hiérarchie des niveaux de densité (voir la \Fig[fig:hartigan], ainsi que la \Fig[fig:hdbscan_eom]).

Évidemment, une première approche aurait pu consister à faire des hypothèses sur la distribution $f$ et supposer, par exemple, un mélange gaussien. Notre modèle eût alors été paramétrique\footnote{De paramètres dans le cas du mélange gaussien : les centres, les matrices de covariances ainsi que les poids associés à chaque gaussienne.}, ce qui nous aurait conduit à la question du choix des `bons' estimateurs pour ces paramètres.
Cette approche, qu'on pourrait qualifier de «~générative~», n'a pas du tout été la nôtre ; la voie que nous avons choisie fut «~discriminative~» (se reporter à l’introduction du chapitre \ref{chap:single-linkage_trois_vues} pour une brève discussion sur ce sujet).

En ne posant aucun \emph{a priori} sur la densité $f$, il ne nous restait guère que l'estimateur $\widehat f_ {K\text{-NN}}$ des $K$ plus proches voisins pour pouvoir identifier la hiérarchie des niveaux de densité.

C'est alors que se produit un ``miracle'' pour $K=1$ : les niveaux, ou \emph{amas}, de forte densité sont \emph{exactement} donnés par un objet très simple : un graphe géométrique. 
Mieux encore : toute l'information hiérarchique -- pour tous les niveaux -- est comprise dans un seul graphe, le plus petit des graphes : l'arbre minimum couvrant sur les données $\mathcal X$.

Ces propriétés font le pont entre le modèle de \nom{Hartigan}, mathématiquement élégant et satisfaisant, à un algorithme heuristique très courant : le Single-Linkage.

Ce constat nous a mis le pied à l'étrier ; la question suivante était naturellement amenée :

\Citation{Si la hiérarchie des niveaux de densité est donnée par un simple graphe lorsque la densité est estimée à l'aide de $K=1$ voisin, qu'en est-il lorsque $K \ge 2$ ?}
La réponse à cette question s'est avérée beaucoup plus riche et dépassant le seul cadre euclidien du modèle de \nom{Hartigan} ; elle constitue l'un des deux piliers fondamentaux de cette thèse :
\Citation{Quand les relations binaires d'un graphe (entre paires de points) se révèlent insuffisantes ou trop instables, il est nécessaire d'intégrer des \textbf{interactions d'ordre supérieur} pour garantir une cohérence robuste.}

Toutes les propriétés du Single-Linkage, tant théoriques qu'algorithmiques, se généralisent naturellement lorsque l'on passe ainsi des graphes aux \emph{hyper}graphes.

Le second pilier de notre thèse concerne l'analyse de cette famille d'algorithmes que l'on peut construire à partir de graphes ou d'hypergraphes. L'on pourrait le résumer ainsi :
\Citation{La théorie de la \textbf{percolation} est la clef de voûte théorique pour comprendre et quantifier les performances de tels algorithmes.}

Là encore, notre analyse ne s'est pas limitée au cadre euclidien, mais a été poussée jusqu'à des données en entrée autrement plus structurées que de simples nuages de points : des graphes (dont les arêtes portent des poids et les nœuds des attributs) pour lesquels il s'agit de retrouver les «~communautés~», ou des images dont on souhaite extraire automatiquement un réseau de fissures.

Nous donnons ci-dessous un bilan des contributions plus détaillé et faisons le point, chapitre par chapitre, sur les questions ouvertes et autres perspectives qui pourraient suivre nos travaux.

\section{Bilan des contributions}

Notre cheminement s'est articulé autour de quatre axes majeurs :
\begin{enumerate}
    \item \textbf{Une vue unifiée pour le clustering hiérarchique.} Tout l'intérêt de la première partie est d'avoir proposé un cadre unifié et cohérent pour le Single-Linkage. Nous nous sommes concentrés sur les trois points de vue qui nous intéressaient pour généraliser naturellement cet algorithme, à savoir une vue qui peut être :
    \begin{enumerate}
        \item \textbf{Heuristique.} Le Single-Linkage est souvent décrit par une heuristique qui revient à construire l'arbre minimum couvrant sur les données. C'est aussi de cette façon qu'il est mis en œuvre. Pour autant, il possède une grande richesse théorique si l'on ne le restreint pas à ce cadre heuristique.
        \item \textbf{Geométrique.} Définir le Single-Linkage au moyen de (hyper)graphes géométriques permet de le relier au phénomène mathématique de la percolation, point de vue indispensable pour analyser ses performances théoriques.
        \item \textbf{Statistique.} C'est le modèle de \nom{Hartigan}, qui explique exactement ce que l'on cherche (les amas de forte densité), et à quels estimateurs l'on recourt pour ce faire (l'estimateur des $K$ plus proches voisins).
    \end{enumerate}
    Fait révélateur qui montre l'importance -- et l'omniprésence -- du Single-Linkage : on retrouve cet algorithme tout aussi bien en informatique théorique, si l'on souhaite proposer une axiomatique minimale de ce que devrait être un «~bon~» clustering\footnote{Ce point de vue ne constitue qu'une digression au cours de nos recherches, celui-ci ne se prêtant pas aussi naturellement à la généralisation que les autres approches. Le choix d'une axiomatique reste, en dernière analyse, toujours assez arbitraire.}.
    Ce dernier point permet de voir le problème de clustering comme une application des espaces métriques vers les espaces ultramétriques. Ces espaces ultramétriques sont d'ailleurs des espaces de recherche privilégiés : nous montrons sur des problèmes industriels complexes comment y injecter simplement des \emph{a priori} géométriques et obtenir des réponses rapides et satisfaisantes, et ce, sans entraînement. Les deux problèmes qui nous servent à illustrer la conclusion de la première partie sont la détection d'anomalies 3D LiDAR et la segmentation panoptique 4D.
    \item \textbf{\HGP-Clusterer : la généralisation par l'ordre supérieur.} De nombreux algorithmes ont été proposés pour rendre davantage robuste le Single-Linkage, le plus universellement adopté aujourd'hui est \HDBSCAN. Ne recourant pas à des interactions d'ordre supérieur (ils continuent de travailler sur des graphes), ces algorithmes créent une asymétrie entre les contraintes fortes qu'ils imposent aux nœuds, tandis qu'ils ne peuvent mettre que des contraintes permissives sur la connexion par arêtes.  Notre algorithme, \HGP-Clusterer, lui, identifie exactement les amas de forte densité des $K$ plus proches voisins pour n'importe quel entier $K$. Pour ce faire, nous avons substitué aux graphes géométriques des hypergraphes géométriques, à savoir : les complexes de \v Cech, et aux composantes connexes les $K$-polyèdres. Sur le plan algorithmique, nous avons prouvé que les fusions utiles (les simplexes $K$-séparants) vérifient une condition de vacuité stricte (être un simplexe de \nom{Gabriel}) et peuvent être extraites efficacement à partir de la mosaïque de \nom{Delaunay} d'ordre supérieur. Nous montrons d'ailleurs sur des données synthétiques 2D et un jeu de données d'huiles d'olive italiennes (représentées par des vecteurs de $\mathbb R^8$) que prendre $K=2$ suffit déjà à obtenir de très bons résultats.
    \item \textbf{La percolation comme mesure de performance théorique.} La limite historique du Single-Linkage en dimension $p \ge 2$ (l'effet de chaînage) n'est autre qu'un phénomène de percolation continue : le bruit de fond percole et fusionne différents amas avant que ceux-ci n'aient pu être intégralement récupérés. Le Single-Linkage est consistant uniquement en dimension $p=1$ justement parce que le phénomène de percolation ne s'y produit pas. Nous avons formalisé ce mécanisme en introduisant la \emph{vitesse de percolation}. Nous avons ainsi pu quantifier mathématiquement à quel point le Robust Single-Linkage et (H)\DBSCAN échouent à recouvrer correctement les amas de forte densité (leur connectivité, lâche, provoque une percolation prématurée et finalement pas assez rapide). À l'inverse, l'exigence d'interactions d'ordre supérieur ($K$-polyèdres) retarde l'apparition de la composante géante dans le bruit, faisant correspondre exactement l'algorithme avec le modèle statistique de \nom{Hartigan} et permettant de récupérer une fraction théorique bien supérieure des amas.
    \item \textbf{Passage à des données davantage structurées : le \emph{leitmotiv} percolation/ordre supérieur demeure.} La troisième partie du manuscrit a montré que l'on peut aller au-delà du paradigme euclidien et retrouver les mêmes phénomènes à l'œuvre (robustification grâce aux interactions d'ordre supérieur et percolation pour l'analyse théorique):
    \begin{itemize}
        \item \textbf{Sur les graphes signés (détection de communautés).} Nous avons prouvé, dans un cadre bayésien soigneusement établi, que le recouvrement partiel des communautés est strictement conditionné par la percolation d'une dynamique de clusters. En généralisant la dynamique de \SW{} avec des interactions d'ordre supérieur (dynamique triangulaire), nous réduisons la frustration géométrique, retardons la percolation parasite et obtenons des bornes théoriques de recouvrement strictement plus fines. % (notamment sur la grille triangulaire pour le modèle de \SB{}).
        \item \textbf{En traitement d'image (extraction de réseaux de fissures).} Face au manque de généralisation de l'apprentissage profond lors de changements de domaine (\emph{domain shifts}), nous avons introduit le graphe de \nom{Frangi} multimodal. L'imposition d'une connexité d'ordre supérieur ($K=2$) agit comme un filtre qui empêche le bruit local de percoler en fausses branches, offrant une extraction robuste des fissures dans des milieux fortement granulaires.
    \end{itemize}
\end{enumerate}

\section{Ouvertures et perspectives de recherche}

Les travaux présentés dans cette thèse ouvrent plusieurs directions de recherche. Nous en proposons quelques-unes en reprenant le fil de la thèse.

%Revues par thèmes. Distinguer théoriques / pratiques 

%* **Champs gaussiens et percolation asymptotique.** L'étude de la vitesse de percolation à haut rappel lorsque $K \to +\infty$ n'a été qu'esquissée. Une analyse rigoureuse des obstructions topologiques des K-polyèdres face aux excursions gaussiennes (via l'espace de Cameron-Martin et le théorème de Schilder) permettrait de clore mathématiquement la preuve de la supériorité asymptotique de HGP-Clusterer.
\begin{itemize}
    \item \textbf{Chapitre \ref{chap:axiomatique}.} L'approche axiomatique n'est pas celle que nous avons privilégiée, pourtant elle pourrait servir à chercher et identifier les bons espaces de recherche (\emph{clustering domain}). Par exemple, les métriques d'arbre pourraient être un bon objet à regarder, légèrement plus général que les ultramétriques (qui sont en correspondance avec un type particulier d'arbre : les dendrogrammes). De plus, les outils rencontrés dans ce chapitre, comme la distance de \nom{Gromov--Hausdorff}, seraient un autre moyen de quantifier la stabilité de notre \HGP-Clusterer.
    \item \textbf{Chapitre \ref{chap:limites_single_linkage} \& \ref{sec:chap_percolation}.} \HGP-Clusterer n'est pas plus consistant (au sens de \nom{Hartigan}) que le Single-Linkage. Il est seulement \emph{fractionnellement} consistant (et cette fraction se mesure grâce à la théorie de la percolation). La question d'avoir un algorithme qui soit consistant en toute dimension $p \ge 2$ reste ouverte. Pour l'obtenir, il serait nécessaire d'introduire des notions de composantes connexes encore plus restrictives.
    \item \textbf{Chapitre \ref{sec:chap_percolation}.} Le champ gaussien introduit au cours du chapitre, obtenu comme limite du \emph{shot noise} centré réduit, possède certainement de bonnes propriétés mathématiques qui le rendent intéressant pour des problèmes théoriques comme celui du \emph{sphere packing} en grande dimension, pour retrouver le dernier facteur 2 manquant à \cite{Campos2023SpherePacking} par rapport à la conjecture de \nom{Parisi} \& \nom{Zamponi}. De manière générale, notre travail apporte des perspectives à toute démonstration faisant intervenir des graphes géométriques (au sens de \nom{Penrose}).  
    \item \textbf{Chapitre \ref{chap:generalisation_mst}.} \textbf{Réduction de dimension d'ordre supérieur.} La mosaïque de Delaunay d'ordre $K$ contient une information structurelle autrement plus riche qu'un simple graphe de voisinages utilisé par des méthodes comme UMAP. Construire un algorithme de réduction de dimension préservant la géométrie dictée par cet objet (ou une «~bonne~» approximation) constitue un prolongement naturel de notre étude et notre piste privilégiée.
    %\item \textcolor{red}{Parler du champ d'excursion gaussien et de problèmes théoriques ``continuisés'' par rapport aux graphes géométriques ? (\emph{sphere packing} en grande dimension.)}
    \item \textbf{Chapitre \ref{chap:bayes_clusters_graphes}.} \textbf{Au-delà du couplage à une réplique.} Bien que les interactions d'ordre supérieur aident fondamentalement à repousser les limites de la frustration pour obtenir une meilleure fraction de recouvrement, elles ne permettent pas d'atteindre le véritable seuil critique de recouvrement faible. Notre cadre bayésien ne construit en effet qu'une seule réplique de la vérité terrain $\Sigma$ par couplage avec une dynamique de type \SW. Disposer de deux répliques conditionnellement indépendantes permettrait d'obtenir des égalités sur le seuil de recouvrement (et non plus seulement des bornes), fournissant ainsi une condition à la fois nécessaire et suffisante.
    \item \textbf{Chapitre \ref{chap:frangi_fissures}.} Il nous paraît de plus en en plus évident que la méthode ``\nom{Frangi}-(hyper)graphe'' que nous avons développée trouverait toute son utilité au sein d'un logiciel d'aide à l'annotation experte de vérité terrain pour des images de réseaux de fissures. L'avantage de notre méthode est qu'elle rend en sortie un (hyper)graphe,  objet aisément ``éditable''. Le point fort de notre méthode étant que nous avons montré qu'elle peut être appliquée à une très grande variété d'images, et ce sans apprentissage, en réglant simplement de façon \emph{ad hoc} quelques paramètres (comme la taille des fissures que l'expert souhaiterait annoter).
\end{itemize}
 et avant \backmatter, par exemple.
% Version refondue : cartographie en niveaux de gris.

\appendix
\chapter{Cartographie des publications}
\cleartorecto
\begin{sidewaysfigure}[p]
\centering
\resizebox{1.00\textheight}{!}{%
\begin{tikzpicture}[x=1cm,y=1cm]

% -----------------------------------------------------------------------------
% Charte locale : uniquement des nuances de gris.
% -----------------------------------------------------------------------------
\colorlet{CartBlack}{black!88}
\colorlet{CartDark}{black!72}
\colorlet{CartMid}{black!46}
\colorlet{CartLight}{black!14}
\colorlet{CartPale}{black!5}
\colorlet{CartVeryPale}{black!2}

\tikzset{
  cartTitle/.style={
    font=\Large\bfseries\scshape,
    text=CartBlack,
    align=center
  },
  cartSubtitle/.style={
    font=\small,
    text=CartDark,
    align=center
  },
  colTitle/.style={
    font=\bfseries\small\scshape,
    text=CartDark,
    anchor=west
  },
  manuscript/.style={
    draw=CartMid,
    fill=CartVeryPale,
    rounded corners=1.8mm,
    line width=.65pt,
    inner sep=2.0mm,
    align=left,
    font=\scriptsize
  },
  axisbox/.style={
    draw=CartBlack,
    fill=CartPale,
    rounded corners=2.0mm,
    line width=.9pt,
    inner sep=2.1mm,
    align=left,
    font=\scriptsize
  },
  paper/.style={
    draw=CartMid,
    fill=white,
    rounded corners=1.8mm,
    line width=.7pt,
    inner sep=1.9mm,
    align=left,
    font=\scriptsize
  },
  exactlink/.style={
    -{Latex[length=2.2mm,width=1.65mm]},
    draw=CartBlack,
    line width=1.25pt,
    rounded corners=2mm
  },
  progresslink/.style={
    -{Latex[length=2.2mm,width=1.65mm]},
    draw=CartBlack,
    line width=1.25pt,
    dashed,
    dash pattern=on 3.2pt off 2.2pt,
    rounded corners=2mm
  },
  structlink/.style={
    -{Latex[length=1.9mm,width=1.4mm]},
    draw=CartDark,
    line width=1.0pt,
    rounded corners=2mm
  },
  note/.style={
    font=\scriptsize,
    text=CartDark,
    align=left
  }
}

\newcommand{\tightenum}{%
  \setlength{\leftmargini}{4.4mm}%
  \setlength{\itemsep}{0.25mm}%
  \setlength{\parsep}{0pt}%
  \setlength{\topsep}{0.35mm}%
}

% -----------------------------------------------------------------------------
% En-tête
% -----------------------------------------------------------------------------
\node[cartTitle] at (14.5,17.25)
  {Cartographie manuscrit -- axes thématiques -- publications};
\node[cartSubtitle] at (14.5,16.70)
  {lecture de gauche à droite : structure du manuscrit, quatre axes avec sous-axes, puis articles publiés, soumis ou en rédaction};

\node[colTitle] at (0.25,16.05) {Structure du manuscrit};
\node[colTitle] at (8.65,16.05) {Axes principaux};
\node[colTitle] at (19.15,16.05) {Articles et travaux associés};

\draw[CartLight, line width=.6pt] (8.10,0.45) -- (8.10,15.78);
\draw[CartLight, line width=.6pt] (18.55,0.45) -- (18.55,15.78);

% -----------------------------------------------------------------------------
% Colonne gauche : structure du manuscrit
% -----------------------------------------------------------------------------
\node[manuscript, anchor=west] (msFond) at (0.25,13.95) {%
\begin{minipage}{7.25cm}
\textbf{Introduction + Partie I}\\[-.3mm]
\textit{Fondations du Single-Linkage}\\[-.4mm]
Ch.~1--5 : trois lectures du Single-Linkage, point de vue axiomatique, limites du chaînage, Robust SL, DBSCAN/HDBSCAN, arbre hiérarchique comme espace de résolution.
\end{minipage}};

\node[manuscript, anchor=west] (msHGP) at (0.25,10.55) {%
\begin{minipage}{7.25cm}
\textbf{Partie II}\\[-.3mm]
\textit{HGP-Clusterer et percolation}\\[-.4mm]
Ch.~6--9 : complexe de \nom[Cech]{\v Cech}, $K$-polyèdres, amas $K$-NN, vitesse de percolation, simplexes $K$-séparants, \nom{Gabriel}, \nom{Delaunay} d'ordre $K$, pratique $K=2$.
\end{minipage}};

\node[manuscript, anchor=west] (msImage) at (0.25,6.35) {%
\begin{minipage}{7.25cm}
\textbf{Partie III, ch.~10 et ch.~12}\\[-.3mm]
\textit{Données image et signal spatial}\\[-.4mm]
Graphe de \nom{Frangi}, hypergraphe de \nom{Frangi}, extraction de réseaux de fissures, multimodalité visible--profondeur, robustesse au bruit et aux changements de domaine.
\end{minipage}};

\node[manuscript, anchor=west] (msGraph) at (0.25,2.65) {%
\begin{minipage}{7.25cm}
\textbf{Partie III, ch.~10 et ch.~11}\\[-.3mm]
\textit{Graphes signés et communautés}\\[-.4mm]
Cadre bayésien, dynamiques de clusters, frustration, interactions triangulaires, percolation comme obstruction ou garantie de recouvrement.
\end{minipage}};

% -----------------------------------------------------------------------------
% Colonne centrale : axes et sous-axes avec enumerate
% -----------------------------------------------------------------------------
\node[axisbox, anchor=west] (axHO) at (8.65,13.60) {%
\begin{minipage}{8.85cm}
\textbf{Axe 1 -- Interactions d'ordre supérieur}\\[-1mm]
\begin{enumerate}\tightenum
  \item Connexité $K$ : $K$-polyèdres, simplexes de \nom[Cech]{\v Cech} et hypergraphes.
  \item Calcul : $K$-arbre minimum, graphes de \nom{Gabriel} et mosaïque de \nom{Delaunay} d'ordre $K$.
  \item Transferts : interactions triangulaires dans les graphes signés et dans les images.
\end{enumerate}
\end{minipage}};

\node[axisbox, anchor=west] (axPerc) at (8.65,9.90) {%
\begin{minipage}{8.85cm}
\textbf{Axe 2 -- Percolation}\\[-1mm]
\begin{enumerate}\tightenum
  \item Lecture géométrique : composantes géantes dans les graphes ou complexes persistants.
  \item Mesure de performance : vitesse de percolation et fraction asymptotiquement récupérable.
  \item Rôle transversal : filtrage du bruit, séparation de niveaux de densité, obstruction au recouvrement.
\end{enumerate}
\end{minipage}};

\node[axisbox, anchor=west] (axSignal) at (8.65,6.05) {%
\begin{minipage}{8.85cm}
\textbf{Axe 3 -- Traitement du signal et des images}\\[-1mm]
\begin{enumerate}\tightenum
  \item Passage pixel $\rightarrow$ arête : graphe de \nom{Frangi} et réponse par paires de pixels.
  \item Extension multimodale : intensité, profondeur, visible--thermique.
  \item Sortie structurelle : extraction d'un réseau linéique plutôt qu'une simple segmentation binaire.
\end{enumerate}
\end{minipage}};

\node[axisbox, anchor=west] (axComm) at (8.65,2.55) {%
\begin{minipage}{8.85cm}
\textbf{Axe 4 -- Détection de communautés sur graphes}\\[-1mm]
\begin{enumerate}\tightenum
  \item Modèles : graphes euclidiens signés, postérieure bayésienne et énergie de \nom{Gibbs}.
  \item Algorithmes : généralisation de \nom{Swendsen--Wang} par liens d'ordre supérieur.
  \item Garanties : recouvrement partiel, presque parfait ou parfait, avec bornes de percolation.
\end{enumerate}
\end{minipage}};

% -----------------------------------------------------------------------------
% Colonne droite : articles. Les flèches portent le statut de correspondance.
% -----------------------------------------------------------------------------
\node[paper, anchor=west] (pANS) at (19.15,14.65) {%
\begin{minipage}{9.45cm}
\textbf{Journal central HGP}\\[-.35mm]
\emph{Generalization of single-linkage with higher-order interactions} : $K$-polyèdres, amas $K$-NN, $K$-arbre et optimisations $K=2$ \cite{BibiAppliedNetworkScience}.
\end{minipage}};

\node[paper, anchor=west] (pHGPConf) at (19.15,12.85) {%
\begin{minipage}{9.45cm}
\textbf{Versions conférence intégrées puis dépassées}\\[-.35mm]
Hypergraphes pour le clustering densité \cite{BibiLyon}; hypergraphes, percolation et clustering hiérarchique \cite{BibiIstambul}; mesure théorique des performances \cite{BibiVenise}.
\end{minipage}};

\node[paper, anchor=west] (pPerco) at (19.15,9.95) {%
\begin{minipage}{9.45cm}
\textbf{Percolation comme fil directeur}\\[-.35mm]
Le chapitre de percolation reprend \cite{BibiIstambul} mais le pousse au cadre continu, aux vitesses critiques et aux limites gaussiennes; \cite{BibiVenise} fournit l'indice de performance (la \emph{vitesse de percolation}).
\end{minipage}};

\node[paper, anchor=west] (pPrec) at (19.15,8.05) {%
\begin{minipage}{9.45cm}
\textbf{Préfigurations}\\[-.35mm]
Extraction de filaments de galaxies par graphes géométriques et percolation \cite{BibiMenton}; complexes simpliciaux et stylométrie comme antécédent méthodologique \cite{BibiMasterBB}.
\end{minipage}};

\node[paper, anchor=west] (pGRETSI) at (19.15,5.75) {%
\begin{minipage}{9.45cm}
\textbf{Première version image}\\[-.35mm]
Graphe de \nom{Frangi} pour les fissures des Vaches Noires, sans données annotées pour la méthode géométrique \cite{BibiStrasbourg}. % ; version EUSIPCO 2025 soumise \cite{BibiPalerme}
%; article de journal en préparation \cite{BibiISPRS}.
\end{minipage}};

\node[paper, anchor=west] (pISPRS) at (19.15,4.05) {%
\begin{minipage}{9.45cm}
\textbf{Versions en préparation / soumise}\\[-.35mm]
Extension multimodale, accélération GPU, ordre supérieur et robustesse bruit/domaine : journal ISPRS en préparation \cite{BibiISPRS} ; EUVIP'26 soumis \cite{BibiLuxembourg}.
\end{minipage}};

\node[paper, anchor=west] (pAsilomar) at (19.15,2.20) {%
\begin{minipage}{9.45cm}
\textbf{Version conférence communautés}\\[-.35mm]
Dynamiques de clusters Monte-Carlo d'ordre supérieur pour la détection de communautés sur graphes euclidiens \cite{BibiMonterey}.
\end{minipage}};

\node[paper, anchor=west] (pInfo) at (19.15,0.85) {%
\begin{minipage}{9.45cm}
\textbf{Article journal en cours de rédaction}\\[-.35mm]
Cadre bayésien pour graphes signés, bornes plus fines et chaînes de Markov à interactions d'ordre supérieur \cite{BibiInformationInference}.
\end{minipage}};

% -----------------------------------------------------------------------------
% Flèches : structure -> axes
% -----------------------------------------------------------------------------
\draw[structlink] (msFond.east) -- (axHO.west);
\draw[structlink] (msHGP.east) -- (axHO.west);
\draw[structlink] (msHGP.east) -- (axPerc.west);
\draw[structlink] (msImage.east) -- (axSignal.west);
\draw[structlink] (msGraph.east) -- (axComm.west);

% -----------------------------------------------------------------------------
% Flèches : axes -> articles.
% Pleine : correspondance directe. Pointillée : progrès significatif dans le manuscrit.
% -----------------------------------------------------------------------------
\draw[exactlink] (axHO.east) -- (pANS.west);
\draw[progresslink] (axHO.east) -- (pHGPConf.west);

\draw[progresslink] (axPerc.east) -- (pPerco.west);
\draw[progresslink] (axPerc.east) -- (pPrec.west);

\draw[progresslink] (axSignal.east) -- (pGRETSI.west);
\draw[exactlink] (axSignal.east) -- (pISPRS.west);

\draw[progresslink] (axComm.east) -- (pAsilomar.west);
\draw[exactlink] (axComm.east) -- (pInfo.west);

% -----------------------------------------------------------------------------
% Légende
% -----------------------------------------------------------------------------
\node[note, anchor=west] at (0.25,0.80) {%
\begin{tabular}{@{}ll@{}}
\tikz\draw[exactlink] (0,0)--(.9,0); & correspondance directe entre manuscrit et article ;\\[1mm]
\tikz\draw[progresslink] (0,0)--(.9,0); & progrès significatif dans le manuscrit par rapport à l'article ;\\[1mm]
\tikz\draw[structlink] (0,0)--(.9,0); & rattachement d'un chapitre à l'un des quatre axes.
\end{tabular}};

% \node[font=\scriptsize\itshape, text=CartMid, anchor=east] at (28.95,0.25)
%   {Nuances de gris uniquement ; flèches épaisses.};

\end{tikzpicture}%
}
\caption[Cartographie manuscrit-publications]{Cartographie des relations entre la structure du manuscrit, les quatre axes thématiques et les publications associées.}
\label{fig:annexe_cartographie_articles}
\end{sidewaysfigure}
